\problemname{Gosedjursmani}
Olivia och Anton är på Liseberg, och de har hittat en kloautomat med $N$ stycken gosedjur, numrerade från $1$ till $N$.
Olivia är väldigt skicklig, 
och vill visa upp sina färdigheter genom att fiska upp Antons favoritdjur.
Den tvådimensionella kloautomaten kan modelleras som ett rutnät av storlek $R \times C$, 
med $R$ rader och $C$ kolumner. 
I den ligger gosedjuren tätt packade, 
och varje ruta hör till ett gosedjur.
Ett gosedjur kan bestå av flera rutor i rutnätet, och varje gosedjur bildar
en sammanhängande figur. 
Med andra ord kan du rita varje gosedjur med en tjock penna utan att lyfta pennan från pappret.
I ett drag kan Olivia välja ett gosedjur som inte täcks av något annat uppifrån, 
och lyfta upp det ur kloautomaten.
Eftersom Olivia inte ännu frågat Anton om vilket hans favoritdjur är, 
vill hon för varje djur veta hur många drag hon sammanlagt behöver göra innan hon kan lyfta upp det djuret.

\noindent

\section*{Indata}
På den första raden står talen $N$, $R$ och $C$ ($1 \leq N \leq R \times C \leq 4 \times 10^{5}$).
Därefter följer $R$ rader, som beskriver rutnätet. 
Den första raden beskriver toppen av kloautomaten, och den sista beskriver botten.
På den $i$:te raden står $C$ heltal $a_{i,1}, a_{i,2}, ..., a_{i,C}$ ($1 <= a_{i,j} <= N$).
Gosedjur nummer $x$ består av alla rutor $i,j$ som uppfyller $a_{i,j} = x$.
Det är garanterat att varje gosedjur består av minst en ruta, och att varje gosedjur bildar en sammanhängande region.

\section*{Utdata}
För varje gosedjur $1, 2, ... N$, skriv ut antalet drag Olivia behöver göra innan hon kan lyfta ut djuret.
Om det inte går, skriv ut $-1$.

\section*{Poängsättning}

\noindent
\begin{tabular}{| l | l | p{12cm} |}
  \hline
  \textbf{Grupp} & \textbf{Poäng} & \textbf{Gränser} \\ \hline
  $1$    & $10$       & $1 \leq N \leq R \times C \leq 20$ \\ \hline
  $2$    & $10$       & $C = 1$ \\ \hline
  $3$    & $10$       & $C \leq 2$ \\ \hline
  $4$    & $10$       & Varje gosedjur har rektangulär form. \\ \hline
  $5$    & $15$       & $N \leq 2000$ \\ \hline
  $6$    & $7$       & $N \leq 30000$ \\ \hline
  $7$    & $38$       & Inga ytterligare begränsningar. \\ \hline
  
\end{tabular}

%% \section*{Förklaring av exempelfall 1}
